\section{Finite-prefix proof script outline}

This chapter outlines how the finite-prefix base is organised as a
proof script in Lean.  The outline describes the structure of the
proof; the repository file is the authoritative source.

\subsection{The goal}

The script proves the conjunction
\begin{equation*}
  \Orbit_7(0)=7\land\Orbit_7(1)=9\land\cdots\land\Orbit_7(16)=34177.
\end{equation*}
In Lean this is
\begin{verbatim}
fullOrbitIter_prefix_expand :
  fullOrbitIter 0 = 7 /\ fullOrbitIter 1 = 9 /\
  ... /\ fullOrbitIter 16 = 34177
\end{verbatim}

\subsection{The unfolding}

The orbit definition unfolds each value:
\begin{equation*}
  \Orbit_7(n+1)=\frac{5\,\Orbit_7(n)+1}{2^{t_n}}.
\end{equation*}
The script rewrites each quotient using the exact division
\begin{equation*}
  5\,\Orbit_7(n)+1=2^{t_n}\,\Orbit_7(n+1).
\end{equation*}

\subsection{Exact division table}

\begin{center}
\scriptsize
\begin{tabular}{@{}rrrrr@{}}
\toprule
$n$ & $\Orbit_7(n)$ & $5\Orbit_7(n)+1$ & $t_n$ & quotient \\
\midrule
0 & 7 & 36 & 2 & 9 \\
1 & 9 & 46 & 1 & 23 \\
2 & 23 & 116 & 2 & 29 \\
3 & 29 & 146 & 1 & 73 \\
4 & 73 & 366 & 1 & 183 \\
5 & 183 & 916 & 2 & 229 \\
6 & 229 & 1146 & 1 & 573 \\
7 & 573 & 2866 & 1 & 1433 \\
8 & 1433 & 7166 & 1 & 3583 \\
9 & 3583 & 17916 & 2 & 4479 \\
10 & 4479 & 22396 & 2 & 5599 \\
11 & 5599 & 27996 & 2 & 6999 \\
12 & 6999 & 34996 & 2 & 8749 \\
13 & 8749 & 43746 & 1 & 21873 \\
14 & 21873 & 109366 & 1 & 54683 \\
15 & 54683 & 273416 & 3 & 34177 \\
\bottomrule
\end{tabular}
\end{center}

\subsection{Weight facts}

The script derives the weight word
\begin{equation*}
  2,1,2,1,1,2,1,1,1,2,2,2,2,1,1,3
\end{equation*}
and then proves the first-big-step conjunction:
\begin{verbatim}
theorem fullOrbit_first_t_ge3_is_exactly_3 :
    (forall n : Nat, n < 15 ->
      twoValuation (5 * fullOrbitIter n + 1) <= 2) /\
    twoValuation (5 * fullOrbitIter 15 + 1) = 3
\end{verbatim}

\subsection{Uniqueness}

The uniqueness lemma is proved by case analysis on $n$:
\begin{itemize}[leftmargin=*]
\item for $n<15$, the small-weight facts contradict the big-weight
  hypothesis;
\item for $n=15$, the table gives weight $3$;
\item for $n>15$, the depth-$15$ step would already be first.
\end{itemize}

\subsection{No \texttt{native\_decide}}

The finite base deliberately avoids \texttt{native\_decide}.  Every
quotient is rewritten from an explicit equality, so the proof stays
inside the kernel's reduction and rewriting machinery.  The audit
searches the main-chain files for \texttt{native\_decide}; the count
is zero.

\subsection{Verification checklist}

\begin{enumerate}[leftmargin=*]
\item \texttt{fullOrbitIter\_\allowbreak prefix\_\allowbreak expand}
  compiles;
\item \texttt{fullOrbit\_\allowbreak prefix\_\allowbreak step\_\allowbreak
  weights} compiles;
\item \texttt{fullOrbit\_\allowbreak first\_\allowbreak t\_\allowbreak
  ge3\_\allowbreak is\_\allowbreak exactly\_\allowbreak 3} compiles;
\item \texttt{first\_\allowbreak big\_\allowbreak step\_\allowbreak
  unique} compiles;
\item no \texttt{native\_decide} appears in the file.
\end{enumerate}

\subsection{Status}

\begin{center}
\small
\begin{tabular}{@{}p{0.56\textwidth}p{0.36\textwidth}@{}}
\toprule
Item & Status \\
\midrule
explicit expansion & compiled \\
step weights & compiled \\
first-big-step facts & compiled \\
uniqueness & compiled \\
no \texttt{native\_decide} & verified by inspection \\
\bottomrule
\end{tabular}
\end{center}
